% The Sage-to-TikZ Figure Pipeline  (documents the sagemath/ pipeline so every
% figure in the book is reproducible)
\chapter{The Sage-to-TikZ Figure Pipeline}
\label{app:C}

Every figure in this book is drawn with exact coordinates, not freehand. The
geometry --- lattice points, the torus and sphere silhouettes, the elliptic curve
samples, the chord and tangent intersections --- is computed by a set of SageMath
scripts that emit \LaTeX{} macros, which the TikZ figures then reference. This
appendix documents the pipeline so the figures can be regenerated and checked.

\section{Overview}
\label{sec:C-overview}

The pipeline lives in the \texttt{sagemath/} directory (a sibling of
\texttt{book/}). Each script \texttt{NN\_name.sage} computes coordinates and
writes them to \texttt{coords\_name.tex} as \verb|\newcommand| definitions; the
driver \texttt{build\_coords.sh} runs all scripts and concatenates their output
into a single \texttt{coords.tex}. The document \verb|\input|s that file, so a
figure can write, e.g., \verb|\fill (\chordPx,\chordPy)| and get the precise
Sage-computed point.

\section{The Scripts}
\label{sec:C-scripts}

\begin{center}
\footnotesize
\begin{tabular}{@{}lll@{}}
\hline
Script & Output \texttt{coords\_*} & Role\\
\hline
\texttt{01\_lattice}          & \texttt{lattice}        & lattice basis and cell data\\
\texttt{02\_torus\_group\_law}& \texttt{group}          & group-law points $u,v,u{+}v$\\
\texttt{03\_elliptic\_curve}  & \texttt{curve}          & the roots $e_1,e_2,e_3$\\
\texttt{04\_chord\_tangent}   & \texttt{chord}          & chord example $P,Q,R,-R$\\
\texttt{05\_tangent\_double}  & \texttt{tangent}        & doubling example $P,R,2P$\\
\texttt{06\_two\_torsion}     & \texttt{torsion}        & two-torsion + vertical-line points\\
\texttt{07\_branched\_cover}  & \texttt{branch}         & branched-cover data\\
\texttt{08\_torus\_render}    & \texttt{torus}          & torus arcs (precise torus)\\
\texttt{09\_torus\_markers}   & \texttt{torus\_markers} & markers on the torus\\
\texttt{10\_sphere\_render}   & \texttt{sphere}         & sphere arcs (the $\CP^1$ image)\\
\texttt{11\_correspondence}   & \texttt{correspondence} & three-panel markers\\
\texttt{12\_curve\_sample}    & \texttt{curve\_sample}  & arcs of $y^2=4x^3-4x$\\
\hline
\end{tabular}
\end{center}

\noindent The macros are named after their script: e.g.\ \texttt{01\_lattice}
emits \verb|\latwax|, \verb|\latomegaone|; \texttt{04\_chord\_tangent} emits
\verb|\chordPx|, \verb|\chordQx|, \verb|\chordRx|, \verb|\chordSx|;
\texttt{11\_correspondence} emits the \verb|\cor...| family; and the three
\verb|\drawprecise...| macros consume \texttt{coords\_torus}, \texttt{coords\_sphere},
and \texttt{coords\_curve\_sample}.

The running example throughout is \(E : y^2 = 4x^3 - 4x\) (\(g_2 = 4\), \(g_3 = 0\),
roots \(-1, 0, 1\)), with a fixed lattice whose basis is recorded in
\texttt{coords\_lattice.tex}.

\section{The Convenience Macros}
\label{sec:C-macros}

Three painting macros, defined in \texttt{book/preamble.tex}, encapsulate the most
intricate figures so a chapter can drop them into a scope:
\begin{itemize}
\item \verb|\drawprecisetorus| --- the tilted torus (outer/inner silhouettes,
  meridians, latitudes split into visible/hidden arcs), from \texttt{coords\_torus};
\item \verb|\drawprecisesphere| --- the \(\CP^1\) sphere matching the torus style,
  from \texttt{coords\_sphere};
\item \verb|\drawprecisecurve| --- the real locus of \(y^2 = 4x^3 - 4x\) (compact
  oval plus unbounded branch, with the three branch points), from
  \texttt{coords\_curve\_sample}.
\end{itemize}

\section{Building and Consuming the Coordinates}
\label{sec:C-build}

The driver \texttt{sagemath/build\_coords.sh} runs every \texttt{0?\_*.sage} in
order and concatenates the per-script files into \texttt{coords.tex}. The book's
preamble loads
\begin{quote}\ttfamily
\textbackslash input\{\textbackslash sagedir/coords.tex\}\\
\textbackslash input\{\textbackslash sagedir/coords\_torus.tex\}\\
\textbackslash input\{\textbackslash sagedir/coords\_curve\_sample.tex\}
\end{quote}
where \verb|\sagedir| points at \texttt{../sagemath}. The aggregated
\texttt{coords.tex} contains the lattice, group, curve, chord, tangent, torsion,
branch, torus-marker, sphere, correspondence, and curve-sample macros; the precise
torus rendering \texttt{coords\_torus.tex} is generated by
\texttt{08\_torus\_render} but loaded separately.

\section{Regenerating the Figures}
\label{sec:C-regen}

With SageMath installed, regenerate all coordinates and rebuild the book by
\begin{quote}\ttfamily
cd sagemath \&\& ./build\_coords.sh\\
cd ../book \&\& pdflatex main.tex \&\& pdflatex main.tex
\end{quote}
The committed \texttt{coords*.tex} files let the book compile \emph{without} Sage;
they need to be regenerated only if a \texttt{.sage} script changes (for instance to
use a different lattice or curve). If a new figure needs precise geometry, add a
script \texttt{sagemath/NN\_name.sage} that writes \texttt{coords\_name.tex}, list
that file in \texttt{build\_coords.sh}, and \verb|\input| it (or rely on the
aggregated \texttt{coords.tex}). Coordinates should never be hand-typed: keeping
them under Sage control is what makes the figures exact and the book reproducible.
